% Cote 119: offprints, reports, courses, CV and letters (typed and printed part): archive cote 119, site page 247 (PDF page 248).
% typescript. Draft compared: none.
% AI-transcribed by Claude Opus 5.5 (claude-opus-5-5, adaptive thinking, effort high), Message Batches API, batch r119_full (msgbatch_01DwZNm5Y7YsKBWyAN4sFs8w), 2026-11-13; not yet checked by eye.
\scanpage{119/247}
\mspage{7}

cours avancé) avait été empêché de venir, aurait fait son cours devant une
salle vide.

4. Le \uline{maître-enseignant} est un homme qui s'intéresse à la science autant ou
plus par les liens humains qu'elle implique, crée ou développe, que par
l'attrait de la découverte scientifique. Sur le plan de la fonction sociale,
c'est donc bien par le rôle d'enseignant qu'il pourra à présent le plus
naturellement rentabiliser sa vocation naturelle. Le plus souvent, il a des
dispositions pour la recherche, qui seules jusqu'à présent étaient censées
l'habiliter à suivre sa vocation principale. Mais souvent aussi, une fois sa
thèse terminée, qui lui aura demandé (à supposer que lui-même ou son
``patron de thèse'' aient pris au sérieux les critères de valeur traditionnels
pour la thèse) un effort très considérable et unique dans sa vie, il n'ap-%
porte plus guère d'autres contributions originales à sa science. Il arrive
même, faute d'aucun mécanisme de contrôle par ses étudiants ni par ses pairs,
qu'il perde tout contact avec la science vivante et qu'il s'installe confor-%
tablement dans son statut de fonctionnaire, répétant au long des années
jusqu'à l'âge de la retraite, sans y changer une virgule, le même cours qu'il
aura mis au point dans sa jeunesse, quand il ne l'a pas hérité tel quel
de son vénérable prédécesseur. Mais nous sortons là de notre description du
maître-enseignant pour celle, trop connue hélas à tous ceux qui ont passé
par l'Université pour qu'il soit utile de la poursuivre, du \uline{maître-fossile}.
Ce dernier, véritable honte de notre métier, et fléau de l'enseignement
depuis l'école communale jusqu'à l'Université, n'est sans doute que par-%
tiellement responsable de son propre rôle stérilisant, n'étant qu'un sous-%
produit de structures scolaires et universitaires qui rendent possible le
foisonnement de lui et de ses pareils. Le vrai enseignant sait d'instinct
qu'il ne saura remplir valablement son rôle qu'à condition d'un renouvellement
constant dans les sujets enseignés, la façon de les présenter, voire même
dans les méthodes d'enseignement.

Pour un enseignant universitaire, un tel renouvellement n'est sans
doute possible qu'en restant au contact \deleted{avec}\inserted{de} la science vivante. Ce n'est
qu'un tel contact qui gardera à son esprit cette ouverture et cette attitude
critique vis-à-vis de la science qu'il enseigne, qui caractérisent la ``tête

\margin{1247}
